\section{原始関数が存在することと式で書けること}\label{節:非初等原始関数}

$\e^{-x^2}$の原始関数を，これまでの置換や部分積分で求めようとしても，なじみのある有限個の式にはたどり着かない。だからといって，原始関数が存在しないわけではない。微積分の基本定理を使えば
\[
 F(x)=\int_0^x\e^{-t^2}\dd t
\]
はすべての実数で定義され，$F'(x)=\e^{-x^2}$となる。これ自体が一つの原始関数である。

「初等関数で表せない」というときは，数，変数，代数的演算，指数関数，対数関数，三角関数などを有限回組み合わせた通常の初等関数の範囲では表せない，という意味である。$\e^{-x^2}$には初等関数で表された原始関数がないことが知られている。ただし，その不可能性の証明には，本書の範囲を超える理論が必要である。ここではこの事実を採用し，置換をいくつか試して失敗したことを証明の代わりにはしない。

積分で定義した$F$についても，多くのことが分かる。被積分関数が正なので$F$は単調増加である。被積分関数が偶関数なので$F$は奇関数である。さらに
\[
 F''(x)=-2x\e^{-x^2}
\]
から，$x>0$では上に凸，$x<0$では下に凸である。有限区間での値はテイラー公式や数値積分で評価できるし，$x\geq1$では$\e^{-x^2}\leq\e^{-x}$なので，$\int_0^{\infty}\e^{-x^2}\dd x$の収束も比較によって分かる。

積分を含む定義は，計算を途中で諦めた記号ではない。微分，対称性，大小関係，近似，極限を調べるための出発点になる。原始関数の表示が難しいときには，求めたいものが「不定積分の式」なのか，「特定の区間での値」なのか，「誤差付きの近似」なのか，「収束するという事実」なのかを分けて考える。その区別により，使う道具も変わってくる。

